\newpage
\phantomsection
\label{prf:preimage-under-inverse-function}
\LRAProofFor{prop:preimage-under-inverse-function}

\begin{remark*}[Return]
\hyperref[prop:preimage-under-inverse-function]{Return to Proposition}
\end{remark*}

\begin{proposition*}[Preimage Under an Inverse Function]
Let $f:A\to B$ be bijective, and let $S\subseteq A$. Then
\[
(f^{-1})^{-1}(S)=f(S),
\]
where $(f^{-1})^{-1}(S)$ denotes the preimage of $S$ under the inverse
function $f^{-1}:B\to A$.
\end{proposition*}

\begin{proof}[Professional Standard Proof]
\LRAProofBodyStart
For $b\in B$,
\[
b\in (f^{-1})^{-1}(S)
\iff f^{-1}(b)\in S
\iff b=f(a)\text{ for some }a\in S
\iff b\in f(S).
\]
Thus $(f^{-1})^{-1}(S)=f(S)$ by extensionality.
\end{proof}

\begin{proof}[Detailed Learning Proof]
\LRAProofBodyStart
Both sets are subsets of $B$, so it is enough to compare their elements. Let
$b\in B$. By the definition of preimage under the function $f^{-1}:B\to A$,
\[
b\in (f^{-1})^{-1}(S)
\quad\Longleftrightarrow\quad
f^{-1}(b)\in S.
\]
Since $f$ is bijective, $f^{-1}(b)$ is the unique element $a\in A$ such that
$f(a)=b$. Therefore $f^{-1}(b)\in S$ exactly when $b=f(a)$ for some
$a\in S$, which is exactly the condition $b\in f(S)$. Hence the two sets have
the same elements.
\end{proof}

\begin{remark*}[Proof structure]
Unpack the outer preimage notation first, then use the defining property of
the inverse function.
\end{remark*}

\begin{dependencies}
\begin{itemize}
  \item \hyperref[def:image-set]{Image of a Set}
  \item \hyperref[def:preimage]{Preimage}
  \item \hyperref[def:inverse]{Inverse Function}
  \item \hyperref[def:bijective]{Bijective Function}
\end{itemize}
\end{dependencies}

\clearpage
